\chapter{System Design and Architecture}
\label{ch:architecture}

\section{Introduction}
System design translates user and functional requirements into an architectural blueprint. This chapter elucidates the high-level architecture of PRIVACY68, followed by formal Data Flow Diagrams (DFD Levels 0, 1, and 2), Unified Modeling Language (UML) structural and behavioral models, and data persistence schemas.

\section{System Architecture}
PRIVACY68 is architected as a modular, event-driven desktop intelligence framework organized into four decoupled subsystems:
\begin{enumerate}
    \item \textbf{Acoustic Ingestion \& Normalization Layer:} Interfaces directly with PortAudio and sound hardware. Captures 16kHz PCM audio strictly on user click-to-speak triggers, applies vector dynamic peak scaling ($\text{norm} = (\text{audio}/\text{peak}) \times 0.9$), executes lexical autocorrection, and invokes the ASR engine.
    \item \textbf{Intent Routing \& Security Execution Layer:} Evaluates transcribed text. Immediate direct commands (media controls, application launches) are dispatched directly. Sensitive or unknown commands pass through the 3-Tier Security Guard, enforcing confirmation dialogues for destructive operations.
    \item \textbf{Universal Multi-Provider Agentic Brain:} Connects via LiteLLM to remote cloud APIs (Google Gemini, Groq, OpenRouter) or local inference daemons (Ollama, LM Studio). Synthesizes executable Python scripts, captures execution stdout/stderr in a sandboxed subprocess, and executes autonomous self-healing retries upon exception tracebacks.
    \item \textbf{Presentation \& Remote Interaction Layer:} Comprises the Microsoft Edge WebView2 Control Center, transparent Win32 HUD overlay, Windows system tray daemon, and a threaded LAN HTTP server providing mobile web remote control.
\end{enumerate}

\begin{figure}[htbp]
\centering
{\scriptsize
\begin{verbatim}
+-------------------------------------------------------------------------+
|                         USER WORKSTATION INTERFACE                      |
|   +-----------------------+                 +-----------------------+   |
|   |   Push-to-Talk HUD    |                 | WebView2 Dashboard    |   |
|   |  (Click-to-Speak UI)  |                 |  (API & Model Config) |   |
|   +-----------+-----------+                 +-----------+-----------+   |
+---------------|-----------------------------------------|---------------+
                | [Trigger Voice]                         | [Config & Keys]
                v                                         v
+-------------------------------+             +-----------------------+
|  Acoustic Capture & Peak Norm |             | Windows Credential    |
|   norm = (audio/peak) * 0.9   |             | Manager & .env Engine |
+---------------+---------------+             +-----------+-----------+
                | [Normalized 16kHz PCM]                  |
                v                                         |
+-------------------------------+                         |
|   ASR Engine (Google Speech)  |                         |
|   [Sub-100ms Transcription]   |                         |
+---------------+---------------+                         |
                | [Transcribed Command]                   |
                v                                         |
+-------------------------------+                         |
|     Intent & Policy Router    |                         |
|  - Fast Path (App/Media)      |                         |
|  - 3-Tier Security Guard      |<------------------------+
+---------------+---------------+
                | [Agentic Request]
                v
+-------------------------------------------------------------------------+
|                  UNIVERSAL AGENTIC BRAIN (LiteLLM)                      |
|   [Primary: Gemini 2.5] --> [Fallback: Groq LPU] --> [Offline: Ollama]  |
+-------------------------------+-----------------------------------------+
                | [Generated Python Code]
                v
+-------------------------------+             +-----------------------+
| Subprocess Execution Sandbox  |             | Persistent Skill DB   |
|  - 30s Timeout                |------------>|     (skills.json)     |
|  - Self-Healing Error Retries |             +-----------------------+
+-------------------------------+
\end{verbatim}
}
\caption{High-Level Architectural Blueprint of PRIVACY68}
\label{fig:high_level_architecture}
\end{figure}

\section{Data Flow Diagrams (DFD)}

\subsection{DFD Level 0: Context Diagram}
The Level 0 Context Diagram depicts the interaction boundary between external entities (the Desktop User, the Mobile Remote User, the Windows Operating System, and External LLM Providers) and the PRIVACY68 system entity.

\begin{figure}[htbp]
\centering
{\scriptsize
\begin{verbatim}
                  [Desktop User]             [Mobile Remote User]
                        |                             |
     Voice / Click Event|                             | HTTP /api/command
     & Confirmation     |                             | & Status Query
                        v                             v
                 +-------------------------------------------+
                 |                                           |
                 |                 PRIVACY68                 |
                 |          Autonomous Desktop System        |
                 |                                           |
                 +-----+-------------------------------+-----+
                       |                               |
       Process Exec,   |                               | Prompt / Completion
       Window Control, |                               | Payload (TLS 1.3)
       File System I/O v                               v
                 [Windows OS]                 [External LLM Providers]
             (Win32 / PowerShell)           (Gemini, Groq, OpenRouter)
\end{verbatim}
}
\caption{DFD Level 0: Context Diagram}
\label{fig:dfd_level_0}
\end{figure}

\subsection{DFD Level 1: System Decomposition Diagram}
The Level 1 DFD decomposes the system into its five constituent operational processes:
\begin{enumerate}
    \item \textbf{Process 1.0 (Acoustic Ingestion):} Captures audio on push-to-talk click, performs peak scaling, and interfaces with ASR.
    \item \textbf{Process 2.0 (Intent Categorization \& Lexical Parsing):} Applies autocorrection, filters vocal tics, and checks for direct app hooks.
    \item \textbf{Process 3.0 (Security Guard Interception):} Validates action risk level against allowlists and coordinates 15-second confirmation timeouts.
    \item \textbf{Process 4.0 (Agentic LLM Synthesis \& Self-Healing):} Generates Python automation code with multi-provider failover.
    \item \textbf{Process 5.0 (OS Subprocess Sandbox Execution):} Spawns isolated execution worker, monitors standard streams, and persists validated skills.
\end{enumerate}

\begin{figure}[htbp]
\centering
{\scriptsize
\begin{verbatim}
[Audio Stream] 
      |
      v
  +-------+  Norm Audio   +-------+   Clean Text   +-------+
  |  1.0  |-------------->|  2.0  |--------------->|  3.0  |
  | Ingest|               | Parse |                | Guard |
  +-------+               +-------+                +---+---+
                                                       | [Approved / Agentic]
                                                       v
[LLM APIs] <======================================+-------+
                                                  |  4.0  |
                                                  | Brain |
                                                  +---+---+
                                                      | [Python Script]
                                                      v
                                                  +-------+
                                                  |  5.0  |====> [Windows OS]
                                                  | Exec  |
                                                  +---+---+
                                                      |
                                                      v
                                              [skills.json DB]
\end{verbatim}
}
\caption{DFD Level 1: Functional Decomposition}
\label{fig:dfd_level_1}
\end{figure}

\subsection{DFD Level 2: Agentic Self-Healing Execution Cycle}
Level 2 DFD zooms into Process 4.0 and 5.0, detailing the autonomous error-recovery feedback loop. If the initial code execution produces an exception or non-zero exit status, the error traceback is dynamically packaged into a correction prompt, returned to the active LLM, and re-executed until verified or until max retries (3) are exhausted.

\begin{figure}[htbp]
\centering
{\scriptsize
\begin{verbatim}
+-------------------+
|  Target Directive |
+---------+---------+
          |
          v
+-------------------+
| LLM Code Synth    |<--------------------------------------+
+---------+---------+                                       |
          | [Python Code]                                   |
          v                                                 | [Retry Context:
+-------------------+                                       |  Code + Traceback]
| Spawn Subprocess  |                                       |
+---------+---------+                                       |
          |                                                 |
          v                                                 |
   Execution Check? =====> [Exit Code != 0 / Traceback] ====+
          |
          | [Exit Code == 0]
          v
+-------------------+
| Success & Notify  |
+---------+---------+
          |
          v
+-------------------+
| Cache to Skill DB |
+-------------------+
\end{verbatim}
}
\caption{DFD Level 2: Self-Healing Code Execution Cycle}
\label{fig:dfd_level_2}
\end{figure}

\section{Unified Modeling Language (UML) Models}

\subsection{Use Case Diagram}
The primary actors interacting with PRIVACY68 are the \textbf{Desktop Operator}, the \textbf{Mobile Remote Client}, the \textbf{External AI Provider}, and the \textbf{Operating System Shell}. Core use cases include:
\begin{itemize}
    \item \texttt{UC-01: Click-to-Speak Audio Command}
    \item \texttt{UC-02: Direct System Application Launch}
    \item \texttt{UC-03: Security Confirmation Approval/Rejection}
    \item \texttt{UC-04: Agentic Multi-Step Automation Synthesis}
    \item \texttt{UC-05: Skill Crystallization and Memory Recall}
    \item \texttt{UC-06: Mobile Remote Control via Local Wi-Fi}
    \item \texttt{UC-07: Provider Key and Model Endpoint Management}
\end{itemize}

\subsection{Sequence Diagram: Agentic Self-Healing Execution}
Figure \ref{fig:uml_sequence} details the temporal message exchange between actors during an agentic voice command execution.

\begin{figure}[htbp]
\centering
{\scriptsize
\begin{verbatim}
User           Streamer         Executor        SecurityGuard    AgentBrain       Subprocess
 |                |                |                  |              |                |
 |--Click Speak-->|                |                  |              |                |
 |  "Clean temp"  |                |                  |              |                |
 |---Release----->|                |                  |              |                |
 |                |--Transcribe--->|                  |              |                |
 |                |                |--Check Risk----->|              |                |
 |                |                |<-Tier 3: Confirm-|              |                |
 |<-Prompt Confirm|                |                  |              |                |
 |--"Yes"-------->|                |                  |              |                |
 |                |                |--Confirm Action->|              |                |
 |                |                |--Generate Code----------------->|                |
 |                |                |                  |              |--Run Script--->|
 |                |                |                  |              |<-ExitCode 1----|
 |                |                |                  |              |  [Traceback]   |
 |                |                |                  |              |--Self Heal---->|
 |                |                |                  |              |--Run Patch---->|
 |                |                |                  |              |<-ExitCode 0----|
 |<-Success HUD---|                |<-Action Done--------------------|                |
\end{verbatim}
}
\caption{UML Sequence Diagram: Self-Healing Execution with Security Interception}
\label{fig:uml_sequence}
\end{figure}

\section{Data Persistence and Configuration Schemas}
PRIVACY68 employs a zero-server, distributed local persistence model designed for maximum privacy, portability, and rapid lookups:

\subsection{Windows Credential Manager Storage}
Sensitive API credentials, specifically fine-grained GitHub Personal Access Tokens (PATs), are stored securely within the Windows Credential Manager via the \texttt{keyring} API under the service namespace \texttt{PRIVACY68/GitHub}. Plaintext tokens are never written to disk.

\subsection{Environment and Model Configuration (\texttt{.env})}
Runtime provider keys and dynamic endpoints are persisted in an application-level \texttt{.env} file loaded via \texttt{python-dotenv}. The schema accommodates:
\begin{itemize}
    \item \texttt{AGENT\_BRAIN\_MODEL}: Configured LLM identifier (e.g., \texttt{gemini/gemini-2.5-flash}, \texttt{groq/llama-3.3-70b-versatile}, \texttt{ollama/qwen2.5:0.5b}).
    \item \texttt{GEMINI\_API\_KEY}: Google AI Studio developer credential.
    \item \texttt{GROQ\_API\_KEY}: Groq Cloud API developer credential.
    \item \texttt{OPENROUTER\_API\_KEY}: OpenRouter multi-provider API credential.
    \item \texttt{OLLAMA\_BASE\_URL}: Custom local/remote endpoint for Ollama daemons.
    \item \texttt{OPENAI\_API\_BASE}: Custom base URL for LM Studio / LocalAI servers.
\end{itemize}

\subsection{Autonomous Skill Database Schema (\texttt{skills.json})}
When an agentic code routine executes successfully, it is stored in \texttt{skills.json} to bypass future LLM inference latency. Each entry adheres to the JSON schema:
\begin{verbatim}
{
  "skill_id": "clean_temp_files",
  "trigger_phrases": [
    "clean temporary files",
    "purge temp directory",
    "clear my temp cache"
  ],
  "keywords": ["clean", "temp", "temporary", "files", "cache"],
  "code": "import os, shutil, tempfile\n...",
  "risk_level": "SENSITIVE",
  "execution_count": 14,
  "last_executed": "2026-11-05T22:15:30Z"
}
\end{verbatim}
This architecture guarantees instantaneous sub-10ms execution for repeated tasks.

